\medskip

\begin{enumerate}
\item \textbf{Affirmation :} La fonction $f$ définie sur $\R$ par $f(x) = \text{e}^x - x$ est convexe.

$f$ est définie et dérivable sur $\R$ comme somme de fonctions dérivables sur $\R$.

On a $f'(x) = \text{e}^x - 1$, puis $f''(x) = \text{e}^x$.

On sait que $\text{e}^x > 0$ quel que soit $x$, donc $f$ est convexe sur $\R$. l'affirmation est vraie.

\item \textbf{Affirmation :}L'équation $\left(2\text{e}^x - 6\right)\left(\text{e}^x + 2\right) = 0$ admet $\ln (3)$ comme unique solution dans $\R$.

Cette équation se traduit par :

\[\left\{\begin{array}{l c l}
2\text{e}^x - 6&=&0\quad \text{ou}\\
\text{e}^x + 2&=&0
\end{array}\right.\]

$\bullet~~$ On sait (voir ci-dessus) que $\text{e}^x > 0 \Rightarrow \text{e}^x + 2  > 2 > 0 $. Donc la deuxième équation n'a pas de solution.

$\bullet~~$$2\text{e}^x - 6=0 \iff 2\text{e}^x =6 \iff \text{e}^x = 3 \iff x = \ln 3$, par croissance de la fonction logarithme népérien : c'est la seule solution donc l'affirmation est vraie.

\item \textbf{Affirmation :}\[\displaystyle\lim_{x \to + \infty} \dfrac{\text{e}^{2x} - 1}{\text{e}^x - x} = 0.\]

On a $\dfrac{\text{e}^{2x} - 1}{\text{e}^x - x} = \dfrac{\text{e}^{2x}\left(1 - \text{e}^{-2x} \right)}{\text{e}^{x}\left(1 - x\text{e}^{-x} \right)} = \text{e}^{x}\dfrac{\left(1 - \text{e}^{-2x} \right)}{\left(1 - x\text{e}^{-x} \right)}$.

$\bullet~~$Au numérateur on sait que $\displaystyle\lim_{x \to + \infty}\text{e}^{-2x} = 0$, donc $\displaystyle\lim_{x \to + \infty}1 - \text{e}^{-2x} = 1$ ;

$\bullet~~$Au dénominateur $x\text{e}^{-x} = \dfrac{x}{\text{e}^x}$ et l'on sait que $\displaystyle\lim_{x \to + \infty}\dfrac{x}{\text{e}^x} = 0$, donc $\displaystyle\lim_{x \to + \infty}1 - x\text{e}^{-x} = 1$.

Le quotient a pour limite 1 et comme $\displaystyle\lim_{x \to + \infty}\text{e}^{x} = + \infty$, on a finalement $\displaystyle\lim_{x \to + \infty} \dfrac{\text{e}^{2x} - 1}{\text{e}^x - x} = +\infty$ : l'affirmation est fausse.

\item La fonction $F$ est définie et dérivable sur $\R$ et sur cet intervalle :

$F'(x) = 2\text{e}^{3x} + 3(2x + 1)\text{e}^{3x} = \text{e}^{3x}(2 + 6x + 3) = (6x + 5)\text{e}^{3x} = f(x)$.

D'autre part $F(0) = 1\text{e}^{0} + 4 = 1 + 4 = 5$

$F$ est donc la primitive de $f$ sur $\R$ qui prend la valeur 5 quand $x = 0$, l'affirmation est vraie.
\item L'execution donne $\dfrac{1 + 9 + \ldots + 5}{10} = \dfrac{50}{10} = 5$.

L'affirmation est fausse.
\end{enumerate}

\bigskip
